\scanpage{53}
\mspage{51}
% archive sheet number: 52
prendraient un sens \unsure{canonique}\deleted{\unsure{ment}}.

Il y a d'autre part un autre point de vue,
qui peut-être pourrait ``canonifier'' les
cat. $\underline{\mathrm{Ch}}_\infty$ en tant qu'objets de $\mathrm{Cat}\,\mathrm{hot}$
et \inserted{surtout} \deleted{\unsure{de}} $\mathrm{Hot}$ (on passerait à $\mathrm{Cat}\,\mathrm{hot}$ en
appliquant le foncteur $\unsure{\Psi|\psi}$, enveloppe de Thomason-Kan \ldots). En effet, on est intéressé à
construire, en termes de $M,\Sigma$, un ``dérivateur''
\inserted{(ou un prédérivateur)} $\unsure{\mathbb D|\mathscr D}_{M,\Sigma}$ \unsure{surtout} sur le ``domaine'' des catégories
finies. L'étendre à un domaine plus large
est peut-être important dans un ``deuxième
souffle'', mais peut être considéré comme
du luxe dans le premier mouvement \inserted{de} fondation
d'une théorie des dérivateurs. Ceci dit,
on peut se poser des questions du type suivant.

Soit \deleted{$(\mathrm{Hot})_0$}
% slanted
\margin{\uline{NB} $\mathrm{Hot}_0$ est équivalente à une petite catégorie}
\[
\mathrm{Hot}_0\ \unsure{\colon\subset|\subset}\ \mathrm{Hot}
\]
la sous-catégorie pleine de $\mathrm{Hot}$, définie par les
cat. finies -- appelons ces types d'homotopie ``de type
fini'' (ou de \uline{présentation} finie ?). On a un
foncteur évident
